% Ad Atticum 16.13 --- headnote
% ad-atticum-16-13, 10 November 44 BC, Aquini

\textit{Cicero to Atticus, written at Aquinum on 10
November 44 BC --- Perseus dateline \textit{Scr. Aquini iv
Id. Nov. a. 710 (44)}. A small filmic moment of the road:
Cicero has risen before dawn at Sinuessa, come at first
light to the Tirenus bridge at Minturnae --- the
turn-off for the inland road to Arpinum --- and there, in
the half-light when there is no lamp yet but not enough
daybreak to read by, Atticus's courier rides up. ``Caught me
pondering a long voyage'' --- Cicero quotes Odyssey 3.169 of
himself.}

\textit{Two letters from Atticus are produced. The first
(``the most graceful of any --- may I not be well if I
write otherwise than I feel'') sets Cicero's course; the
second urges him on with another Odyssean tag (``past windy
Mimas'') --- which Cicero glosses for the joke: keeping the
Appian Way on his left. He stayed that day at Aquinum, set
out next morning, and sent this letter from the road. After
a lacuna, the practical questions: he would rather be at
Tusculum or somewhere in the suburbs --- should he come
closer? --- write often. The closing answer to Atticus's
own question is laconic: if the parties are evenly matched,
sit still; if not, the trouble will reach us, and we will
act together.}
